% Generated draft: requires mathematical and source review.
\documentclass[11pt]{article}
\usepackage[margin=25mm]{geometry}
\usepackage{fontspec}
\setmainfont{Latin Modern Roman}
\usepackage{amsmath,amssymb,mathtools,unicode-math}
\setmathfont{Latin Modern Math}
\usepackage{enumitem,needspace,xurl,hyperref}
\setlength{\parindent}{0pt}
\setlength{\parskip}{5pt}
\setlength{\emergencystretch}{3em}
\newcommand{\N}{\mathbb{N}}
\newcommand{\Z}{\mathbb{Z}}
\newcommand{\Q}{\mathbb{Q}}
\newcommand{\R}{\mathbb{R}}
\newcommand{\C}{\mathbb{C}}
\newcommand{\F}{\mathbb{F}}
\newcommand{\A}{\mathbb{A}}
\newcommand{\PP}{\mathbb{P}}
\newcommand{\E}{\mathbb{E}}
\newcommand{\eps}{\varepsilon}
\newcommand{\dd}{\,\mathrm{d}}
\begin{document}
\section*{UnsolvedMath ID: 30001407}
\textbf{Source problem code:} catalogue entry
\subsection{Definitions and mathematical statement}
The source asks for a proof or disproof of a genus-three conjectural Hecke trace formula. Its parameters satisfy
\[
a\geq b\geq c\geq 0.
\]
The source does not specify their domain, integrality or parity conditions, or the range of the parameter \(p\). The displayed inequalities include equality cases and \(c=0\); no further boundary restrictions are supplied.

The proposed formula, reproduced with the space on which the trace is taken made explicit, is
\[
\operatorname{tr}\!\left(T(p)\mid S_{a-b,b-c,c+4}\right)
=
e_c(A_3\otimes\mathbb{F}_p,V_{a,b,c})
-e_{3,\mathrm{extra}}(a,b,c),
\]
where the supplied correction expression is
\[
\begin{aligned}
e_{3,\mathrm{extra}}(a,b,c)
={}&-e_c(A_2,V_{a+1,b+1})\\
&-e_{2,\mathrm{extra}}(a+1,b+1)\otimes S[c+2]\\
&+e_c(A_2,V_{a+1,c})\\
&+e_{2,\mathrm{extra}}(a+1,c)\otimes S[b+3]\\
&-e_c(A_2,V_{b,c})\\
&-e_{2,\mathrm{extra}}(b,c)\otimes S[a+4].
\end{aligned}
\]
Here \(T(p)\) is described by the source as a Hecke operator on \(S_{a-b,b-c,c+4}\). The source does not define this space or the normalization of the operator. It also supplies no definitions of \(A_2\), \(A_3\), the objects \(V_{u,v}\) and \(V_{a,b,c}\), the invariant \(e_c\), the genus-two correction \(e_{2,\mathrm{extra}}\), or the symbols \(S[k]\) and \(\otimes\). The subscript in \(e_c\) is retained exactly as supplied.

In particular, the source does not specify the mathematical setting in which the correction expression takes values, the base for the \(A_2\) terms, or how the right-hand side is to be compared with the numerical Hecke trace. If its terms are cohomological or motivic objects, an appropriate realization and trace operation would have to be specified; none is supplied. Thus the display records the proposed expression, but a fully defined, quantified conjecture cannot be recovered from this record alone. No additional hypotheses, missing operations, or definitions are imposed here.
\subsection{Short English statement}
Prove or disprove the displayed genus-three Hecke trace formula for parameters satisfying \(a\geq b\geq c\geq0\), with the six supplied correction terms. The record leaves essential domains, definitions, normalizations, and the interpretation of the right-hand side unresolved, so this is a statement-only draft requiring mathematical review.
\subsection{Sources}
\begin{itemize}
\item Supplied local UnsolvedMath registry record, permanent ID \(30001407\), title \emph{Genus Three Hecke Trace Formula}, exact source problem code \texttt{catalogue entry}; collection: \emph{Oberwolfach Reports Open Problems}. The record gives proposed year 2010 and supplies no named proposer.
\item Supplied report citation: \emph{Moduli Spaces in Algebraic Geometry} (2010), \emph{Oberwolfach Reports}, pp.~55--107, DOI \texttt{10.4171/owr/2010/02}. Report and literature-source link, labeled ``original Oberwolfach report'': \url{https://doi.org/10.4171/owr/2010/02}. The supplied background describes these reports as documenting mathematical workshops and their open-problem discussions. The report has not been independently consulted for this draft.
\item Literature assessment dated exactly \texttt{2026-08-21}: the supplied curation labels the problem open and describes the isolated question as a truncated fragment with no identifiable mathematical assertion or hypotheses, for which no resolution could be verified; it retains the open label under its stated policy. This is supplied dated curation, not an independently verified assessment. Its description is in tension with the explicit formula present in the supplied statement. The background expressly requests independent verification of time-sensitive status and cited claims.
\item The supplied statement status is \texttt{corrected\_verified}. Its accompanying verification note says only that TeX notation was enclosed in math delimiters without changing the mathematical content; it supplies no mathematical verification of the formula.
\end{itemize}
\end{document}
